%% ═══════════════════════════════════════════════════════════════════
\chapter*{Bibliografia e Sitografia}
\addcontentsline{toc}{chapter}{Bibliografia e Sitografia}
%% ═══════════════════════════════════════════════════════════════════

%% ─── PROMPT LLM ─────────────────────────────────────────────────
%% Genera la bibliografia completa in formato accademico italiano
%% (stile Chicago Author-Date o APA, coerente con le linee guida
%% dell'ateneo). Organizza in sezioni:
%%
%% A) MANUALI E MONOGRAFIE (fonti teoriche primarie):
%%    - Artoni R. (ultima ed.) -- Elementi di Scienza delle Finanze,
%%      Il Mulino
%%    - Bosi P. (ultima ed.) -- Corso di Scienza delle Finanze,
%%      Il Mulino
%%    - Fornero E. (2018) -- Chi ha paura delle riforme, Egea, Milano
%%    - Rosen H.S., Gayer T., Rapallini C. (2023) --
%%      Scienza delle Finanze, 6a ed., McGraw-Hill Education, Milano
%%    - Thaler R.H., Sunstein C.R. (2008) -- Nudge: Improving Decisions
%%      About Health, Wealth, and Happiness, Yale University Press
%%
%% B) RAPPORTI ISTITUZIONALI (dati e proiezioni):
%%    - MEF/RGS -- Rapporto n.24 sulle tendenze di medio-lungo periodo
%%      del sistema pensionistico e socio-sanitario (2023)
%%    - UPB -- Rapporto sulla politica di bilancio, giugno 2025
%%    - Banca d'Italia -- Relazione del Governatore sul 2024 (2025)
%%    - OCSE -- Pensions at a Glance (ultima edizione disponibile)
%%    - IMF -- Italy Article IV Consultation (anno più recente)
%%    - COVIP -- Relazione Annuale (anno più recente)
%%    - INPS -- Rapporto Annuale (anno più recente)
%%
%% C) ARTICOLI ACCADEMICI E WORKING PAPER:
%%    - Acemoglu D., Restrepo P. (2020) -- "Robots and Jobs: Evidence
%%      from US Labor Markets", Journal of Political Economy, 128(6)
%%    - Fornero E., Monticone C. (CERP) -- studi su adeguatezza
%%      e pensionamento flessibile
%%    - Raitano M. -- proposta pensione contributiva di garanzia,
%%      working paper CERP/INAPP
%%    - Thaler R.H., Benartzi S. (2004) -- "Save More Tomorrow",
%%      Journal of Political Economy, 112(S1)
%%    - Favero C. (2018) -- "Quando le pensioni generose aumentano
%%      il debito", Lavoce.info, 30 novembre 2018
%%
%% D) FONTI DIGITALI E SITOGRAFIA:
%%    - Lavoce.info -- sezione "Pensioni" (articoli citati nel testo)
%%    - VoxEU.org -- articoli su NDC, ageing e automation
%%    - INPS "La Mia Pensione" -- https://www.inps.it
%%    - RGS -- https://www.rgs.mef.gov.it
%% ────────────────────────────────────────────────────────────────
